\documentclass[12pt,a4paper]{article}

\usepackage[utf8]{inputenc}
\usepackage[T1]{fontenc}
\usepackage[spanish]{babel}
\usepackage{geometry}
\usepackage{graphicx}
\usepackage{float}
\usepackage{array}
\usepackage{booktabs}
\usepackage{amsmath}
\usepackage{xcolor}
\usepackage{hyperref}
\usepackage{enumitem}
\usepackage{caption}
\usepackage{fancyhdr}
\usepackage{seqsplit}
\usepackage{csquotes}
\usepackage{listings}
\usepackage{tikz}
\usepackage{everypage}
\usepackage{tikzpagenodes}
\usetikzlibrary{arrows.meta,positioning,shapes.geometric,calc}

% ------------------------------------------------------------------
% IMÁGENES NECESARIAS EN OVERLEAF (subir a la raíz del proyecto):
%   - Logo.png                    -> logo institucional (ESPE)
%   - DCCO.png                    -> logo del Departamento de Ciencias de la Computación
%   - figura1_topologia.png       -> topología VM en modo NAT
%   - figura2_ps_aux.png          -> salida de ps -aux
%   - figura3_ps_aux_more.png     -> salida de ps -aux | more
%   - figura4_ps_ef.png           -> salida de ps -ef
%   - figura5_ps_u_root.png       -> salida de ps -u root
%   - figura6_ps_forest.png       -> salida de ps -e --forest
%   - figura7_top.png             -> herramienta top
%   - figura8_htop.png            -> herramienta htop
%   - figura9_pstree.png          -> pstree básico
%   - figura10_pstree_p.png       -> pstree -p (con PID)
%   - figura11_pstree_u.png       -> pstree -u root
%   - figura12_pstree_a.png       -> pstree -a (con argumentos)
%   - figura13_ls_proc.png        -> ls /proc
%   - figura14_cpuinfo.png        -> cat /proc/cpuinfo
%   - figura15_lscpu.png          -> lscpu
%   - figura16_mpstat.png         -> mpstat global
%   - figura17_mpstat_all.png     -> mpstat -P all 1
% Todas en la misma carpeta que este .tex (proyecto plano de Overleaf).
% NOTA: figuras 18-22 (atop, btop, gtop, glances, nmon) se trasladaron
%        al documento de Deberes 2P.
% ------------------------------------------------------------------

\let\originaltexttt\texttt
\DeclareRobustCommand\texttt[1]{{\ttfamily\seqsplit{#1}}}

\geometry{left=3cm, right=2.75cm, top=4.5cm, bottom=4cm}
\setlength{\parskip}{0.7em}
\setlength{\parindent}{0pt}
\hypersetup{
    colorlinks=true,
    linkcolor=blue,
    urlcolor=blue
}

\pagestyle{empty}

\definecolor{negro_suave}{rgb}{.15,.15,.15}
\definecolor{blanco_puro}{rgb}{1,1,1}

% ------------------------------------------------------------------
\definecolor{codegray}{rgb}{0.95,0.95,0.95}
\definecolor{codegreen}{rgb}{0.0,0.5,0.0}
\definecolor{codeblue}{rgb}{0.0,0.0,0.7}
\definecolor{codepurple}{rgb}{0.58,0.0,0.82}

\lstdefinestyle{academico}{
    backgroundcolor=\color{codegray},
    basicstyle=\ttfamily\footnotesize,
    keywordstyle=\color{codeblue}\bfseries,
    commentstyle=\color{codegreen}\itshape,
    stringstyle=\color{codepurple},
    numbers=left,
    numberstyle=\tiny\color{gray},
    numbersep=8pt,
    frame=single,
    breaklines=true,
    showstringspaces=false,
    tabsize=4,
    captionpos=b,
    literate=%
        {á}{{\'a}}1 {é}{{\'e}}1 {í}{{\'i}}1 {ó}{{\'o}}1 {ú}{{\'u}}1
        {Á}{{\'A}}1 {É}{{\'E}}1 {Í}{{\'I}}1 {Ó}{{\'O}}1 {Ú}{{\'U}}1
        {ñ}{{\~n}}1 {Ñ}{{\~N}}1 {¿}{{?`}}1 {¡}{{!`}}1
}
\lstset{style=academico}

% ------------------------------------------------------------------
% Portada (página 1) y bandas con logos en el resto de páginas.
% ------------------------------------------------------------------
\AddEverypageHook{%
\tikz[remember picture,overlay]{%
  \ifnum\value{page}>1%
    \filldraw[fill=negro_suave, thick, draw=blanco_puro] (1,-1.6) rectangle (16,-1.3);
    \filldraw[fill=negro_suave, thick, draw=blanco_puro] (1,-1.6-22.8) rectangle (16,-1.3-22.8);
        \node[anchor=east] at (16,-25) {\textcolor{black}{{\footnotesize DCCO - ESPE}}};
        \node[anchor=west] at (1,-25) {\textcolor{black}{{\footnotesize \thepage}}};
        \node at (1.8,0) {\includegraphics[width=2.5cm]{DCCO.png}};
  \else
    \filldraw[fill=negro_suave, thick, draw=blanco_puro] (-0.3,0) rectangle (0.3,-26);
    \filldraw[fill=negro_suave, thick, draw=blanco_puro] (1,-1.6) rectangle (16,-1.3);
    \filldraw[fill=negro_suave, thick, draw=blanco_puro] (1,-1.6-21.8) rectangle (16,-1.3-21.8);
    \node[anchor=east] at (16,-25) {\textit{{\Large ING. WALTER FUERTES}}};
    \node[anchor=east] at (16,-25.8) {\textit{{\Large Departamento de Ciencias de la Computación - ESPE}}};
    \node[anchor=east] at (16,-9)  {\textbf{\Huge Sistemas Operativos}};
    \node[anchor=east] at (16,-11) {\textbf{\Large PR\'ACTICA DE LABORATORIO N.\textsuperscript{o} 1 2P}};
    \node[anchor=east] at (16,-13) {\Large Autor: Blacio Machuca Julio Cesar};
    \node[anchor=east] at (16,-14.2) {\Large NRC: 30770};
    \node[anchor=east] at (16,-15.4) {\Large Fecha: 20 de mayo de 2026};
    \node[anchor=east] at (16,-16.6) {\Large Segunda Unidad: Gestión de Procesos};
    \node[anchor=east] at (16,-22) {\textbf{\LARGE GESTI\'ON DE PROCESOS EN LINUX}};
    \node at (4.2,0.) {\includegraphics[width=6cm]{Logo.png}};
    \node at (14,0.) {\includegraphics[width=2.5cm]{DCCO.png}};
  \fi}
}

\begin{document}

% ------------------------------------------------------------------
\thispagestyle{empty}
\null
\newpage

% ------------------------------------------------------------------
\tableofcontents
\newpage
\listoffigures
\listoftables
\newpage

% ==================================================================
\section{Objetivos de aprendizaje}
% ==================================================================

\begin{enumerate}[leftmargin=*]
    \item El estudiante identifique un proceso, lo cree, lo elimine y sea capaz de gestionar
    procesos en primer y segundo plano en un entorno Linux.
    \item Al finalizar la práctica, el estudiante visualice el rendimiento y el desempeño
    del procesador mediante herramientas de monitoreo del sistema.
    \item El estudiante interprete correctamente las métricas de desempeño reportadas por las
    herramientas de monitoreo \texttt{top} y \texttt{htop}, relacionándolas con el estado interno
    de los procesos expuesto por el núcleo a través del sistema de archivos virtual \texttt{/proc}.
\end{enumerate}

% ==================================================================
\section{Marco Teórico}
% ==================================================================

\subsection{Procesos en Linux}

En los sistemas operativos modernos, y en particular en los sistemas basados en Linux, un
\textbf{proceso} se define como un programa en ejecución, es decir, una instancia activa de un
programa que ocupa recursos del sistema y que posee un contexto de ejecución propio. Cada proceso
recibe del kernel un identificador numérico único denominado PID (\emph{Process Identifier}),
comparable a una cédula de identidad digital, que le permite ser distinguido, monitorizado y
administrado de forma individual (Tanenbaum y Bos, 2015). Un proceso difiere de un programa en
que el programa constituye únicamente una secuencia estática de instrucciones almacenadas en
disco, mientras que el proceso representa la ejecución dinámica de dicho programa junto con su
estado en memoria, los registros del procesador, los archivos abiertos y demás recursos activos.

La gestión de procesos constituye una de las funciones centrales del sistema operativo, pues de
ella depende la correcta distribución del tiempo de procesador, la memoria y los dispositivos de
entrada y salida entre los distintos procesos que compiten por estos recursos. En Linux, esta
función recae en el \emph{kernel}, componente que actúa como núcleo del sistema operativo y que
mantiene en todo momento una tabla de procesos con la información de cada proceso activo.

\subsection{Ciclo de vida de un proceso: estados y transiciones}

Durante su existencia, un proceso en Linux atraviesa cinco estados bien definidos, descritos en
la Tabla~\ref{tab:estados}.

\begin{table}[H]
    \centering
    \caption{Estados de un proceso durante su ciclo de vida.}
    \label{tab:estados}
    \begin{tabular}{@{}>{\bfseries}l c p{9cm}@{}}
        \toprule
        Estado      & Símbolo & Descripción \\
        \midrule
        Nuevo       & N & El proceso se está creando. \\
        Listo       & L & El proceso espera la asignación de la CPU por parte del planificador. \\
        Ejecución   & E & El proceso recibe la CPU y ejecuta sus instrucciones. \\
        Bloqueado   & B / S & El proceso espera la finalización de una operación de E/S. \\
        Terminado   & T & El proceso ha concluido su ejecución y libera todos sus recursos. \\
        \bottomrule
    \end{tabular}
\end{table}

Las transiciones entre estos estados siguen una secuencia lógica: el proceso parte del estado
\emph{Nuevo}, avanza al estado \emph{Listo} mediante la admisión del kernel, pasa a
\emph{Ejecución} cuando el planificador le asigna el procesador y, al concluir, transita al
estado \emph{Terminado}. Si durante la ejecución requiere esperar un evento externo, transita
al estado \emph{Bloqueado} y regresa a \emph{Listo} una vez completada la espera.

\subsection{Bloque de Control de Procesos (PCB)}

El sistema operativo mantiene para cada proceso activo una estructura de datos denominada
\textbf{PCB} (\emph{Process Control Block}), que almacena toda la información necesaria para
gestionar y reanudar el proceso en cualquier momento. Sus campos principales son el PID, el
PPID (\emph{Parent Process Identifier}), el estado del proceso, el contador de programa, los
registros de la CPU, los punteros a memoria, los límites de memoria, la lista de archivos
abiertos, la prioridad y los datos contables de uso de recursos (Silberschatz et al., 2018).

\subsection{Tipos de usuarios en Linux}

Linux distingue tres tipos de usuarios en función de sus privilegios sobre el sistema:

\begin{itemize}[leftmargin=*]
    \item \textbf{Usuario sin privilegios (\texttt{\$}):} Puede ejecutar aplicaciones y acceder
    únicamente a los recursos que le han sido asignados explícitamente.
    \item \textbf{Usuario administrador o \texttt{root} (\texttt{\#}):} Posee control total
    sobre el sistema; puede crear, modificar y eliminar cualquier proceso, archivo o
    configuración del sistema operativo.
    \item \textbf{Usuario \texttt{sudo}:} Permite a un usuario sin privilegios ejecutar comandos
    específicos con permisos de administrador, mediante la autenticación con contraseña, sin
    necesidad de iniciar sesión como \texttt{root}.
\end{itemize}

\subsection{Sistema de archivos \texttt{/proc}}

En Linux, el directorio \texttt{/proc} es un sistema de archivos virtual (\emph{procfs}) que no
existe físicamente en el disco, sino que el kernel genera dinámicamente en memoria. Cada proceso
activo tiene un subdirectorio dentro de \texttt{/proc} con nombre igual a su PID, que contiene
archivos con información detallada sobre ese proceso: uso de memoria, estado, línea de comandos
con la que fue lanzado, descriptores de archivos abiertos, entre otros. Adicionalmente, archivos
como \texttt{/proc/cpuinfo}, \texttt{/proc/meminfo} y \texttt{/proc/stat} proveen información
global sobre el hardware y el estado del sistema, siendo la fuente de datos que utilizan
internamente herramientas como \texttt{top} y \texttt{htop} (Love, 2010).

\subsection{Herramientas de gestión y monitoreo de procesos}

Linux ofrece un conjunto de herramientas de línea de comandos para la gestión y el monitoreo de
procesos, clasificadas en dos grupos. Los comandos de consulta y gestión incluyen \texttt{ps},
que genera una fotografía estática de los procesos activos; \texttt{pgrep}, que busca procesos
por nombre o patrón; \texttt{pidof}, que devuelve el PID dado el nombre exacto del ejecutable;
\texttt{sleep}, que crea un proceso de espera controlada; y \texttt{kill}, que envía señales a
los procesos. Las herramientas de monitoreo dinámico comprenden \texttt{top}, que muestra en
tiempo real la lista de procesos ordenados por consumo de CPU; \texttt{htop}, versión interactiva
y mejorada de \texttt{top}; \texttt{pstree}, que representa la jerarquía de procesos en forma de
árbol; y \texttt{mpstat}, que reporta estadísticas de uso del procesador por núcleo (Kerrisk, 2010).

% ==================================================================
\section{Topología}
% ==================================================================

La práctica se desarrolló en un entorno de virtualización compuesto por una máquina virtual Linux
configurada en modo NAT (\emph{Network Address Translation}). En este modo, la máquina virtual
accede a la red a través de la dirección IP del sistema anfitrión, lo que permite la conectividad
a Internet sin exponer la máquina virtual directamente en la red local. La Figura~\ref{fig:topo}
ilustra esta topología.

\begin{figure}[H]
    \centering
    \includegraphics[width=0.65\textwidth]{figura1_topologia.png}
    \caption{Topología de red: máquina virtual en modo NAT.}
    \label{fig:topo}
\end{figure}

% ==================================================================
\section{Desarrollo}
% ==================================================================

\subsection{Sistema de ayuda en línea}

Antes de proceder con la gestión de procesos, se revisaron los mecanismos de ayuda en línea
disponibles en Linux, resumidos en la Tabla~\ref{tab:ayuda}.

\begin{table}[H]
    \centering
    \caption{Comandos de ayuda en línea disponibles en Linux.}
    \label{tab:ayuda}
    \begin{tabular}{@{}l p{9.5cm}@{}}
        \toprule
        Comando & Descripción \\
        \midrule
        \texttt{man <comando>}  & Muestra el manual completo del comando indicado. \\
        \texttt{whatis <cmd>}   & Proporciona una descripción de una sola línea del comando. \\
        \texttt{info <cmd>}     & Presenta documentación extendida en formato hipertextual. \\
        \texttt{<cmd> --help}   & Muestra un resumen rápido de las opciones del comando. \\
        \bottomrule
    \end{tabular}
\end{table}

\subsection{Comando \texttt{ps} — Estado de los procesos}

El comando \texttt{ps} (\emph{Process Status}) genera una lista estática de los procesos activos
en el sistema en el momento en que se ejecuta. A diferencia de \texttt{top}, no actualiza la
información de forma continua. Sus variantes principales se detallan a continuación.

\subsubsection{\texttt{ps -aux} — Todos los procesos de todos los usuarios}

La opción \texttt{-aux} combinó tres flags: \texttt{a} mostró los procesos de todos los
usuarios, \texttt{u} presentó la salida en formato orientado al usuario con columnas de propietario,
PID, \texttt{\%CPU} y \texttt{\%MEM}, y \texttt{x} incluyó los procesos que no estaban asociados a
una terminal. Esta variante resultó la más empleada para obtener una visión global del sistema.

\begin{lstlisting}[language=bash,caption={Comando ps -aux para listar todos los procesos del sistema.},label={lst:psaux}]
$ ps -aux
\end{lstlisting}

\begin{figure}[H]
    \centering
    \includegraphics[width=0.9\textwidth]{figura2_ps_aux.png}
    \caption[Salida de ps -aux.]{Salida del comando \texttt{ps -aux}, mostrando todos los procesos del sistema.}
    \label{fig:psaux}
\end{figure}

Los campos más relevantes de la salida se describen en la Tabla~\ref{tab:campos_ps}.

\begin{table}[H]
    \centering
    \caption{Campos principales de la salida de ps -aux.}
    \label{tab:campos_ps}
    \begin{tabular}{@{}>{\ttfamily}l p{9.5cm}@{}}
        \toprule
        \textrm{Campo} & Descripción \\
        \midrule
        USER    & Nombre del usuario propietario del proceso. \\
        PID     & Identificador único del proceso. \\
        \%CPU   & Porcentaje de tiempo de CPU consumido. \\
        \%MEM   & Porcentaje de memoria RAM utilizada. \\
        VSZ     & Tamaño del espacio de memoria virtual asignado (en KB). \\
        RSS     & Memoria RAM física en uso (\emph{Resident Set Size}, en KB). \\
        STAT    & Estado actual del proceso (S = dormido, R = en ejecución, etc.). \\
        COMMAND & Nombre o ruta del comando que inició el proceso. \\
        \bottomrule
    \end{tabular}
\end{table}

\subsubsection{\texttt{ps -aux | more} — Paginación de la salida}

Dado que la salida de \texttt{ps -aux} puede ser extensa, se canalizó con \texttt{more} para
permitir su revisión paginada:

\begin{lstlisting}[language=bash,caption={Paginación de la salida de ps -aux con more.},label={lst:psmore}]
$ ps -aux | more
\end{lstlisting}

\begin{figure}[H]
    \centering
    \includegraphics[width=0.9\textwidth]{figura3_ps_aux_more.png}
    \caption[Salida de ps -aux paginada con more.]{Salida de \texttt{ps -aux} paginada con \texttt{more}.}
    \label{fig:psmore}
\end{figure}

\subsubsection{\texttt{ps -ef} — Formato estándar completo}

La opción \texttt{-ef} presentó todos los procesos del sistema (\texttt{-e}) en formato
estándar completo (\texttt{-f}), con las columnas UID, PID, PPID, \texttt{\%CPU}, \texttt{STIME},
\texttt{TTY} y \texttt{CMD}. Esta variante resultó especialmente útil para identificar las
relaciones padre-hijo mediante la columna PPID:

\begin{lstlisting}[language=bash,caption={Comando ps -ef con formato estándar completo.},label={lst:psef}]
$ ps -ef
\end{lstlisting}

\begin{figure}[H]
    \centering
    \includegraphics[width=0.9\textwidth]{figura4_ps_ef.png}
    \caption[Salida de ps -ef.]{Salida del comando \texttt{ps -ef}, mostrando PID, PPID y otras métricas.}
    \label{fig:psef}
\end{figure}

\subsubsection{\texttt{ps -u root} — Procesos del usuario root}

La opción \texttt{-u root} filtró y mostró únicamente los procesos pertenecientes al usuario
\texttt{root}, lo que permitió auditar los servicios del sistema con privilegios de administrador:

\begin{lstlisting}[language=bash,caption={Filtrado de procesos del usuario root con ps -u root.},label={lst:psroot}]
$ ps -u root
\end{lstlisting}

\begin{figure}[H]
    \centering
    \includegraphics[width=0.9\textwidth]{figura5_ps_u_root.png}
    \caption[Salida de ps -u root.]{Procesos del usuario \texttt{root} obtenidos con \texttt{ps -u root}.}
    \label{fig:psroot}
\end{figure}

\subsubsection{\texttt{ps -e --forest} — Árbol jerárquico de procesos}

La opción \texttt{--forest} presentó la lista de todos los procesos (\texttt{-e}) en formato de
árbol ASCII que visualizó la jerarquía de procesos padre-hijo mediante sangría y
caracteres de ramificación:

\begin{lstlisting}[language=bash,caption={Árbol de jerarquía de procesos con ps -e --forest.},label={lst:psforest}]
$ ps -e --forest
\end{lstlisting}

\begin{figure}[H]
    \centering
    \includegraphics[width=0.9\textwidth]{figura6_ps_forest.png}
    \caption[Salida de ps -e --forest.]{Árbol de jerarquía de procesos generado con \texttt{ps -e --forest}.}
    \label{fig:psforest}
\end{figure}

\subsection{Herramientas de monitoreo dinámico}

\subsubsection{\texttt{top} — Monitor de procesos en tiempo real}

La herramienta \texttt{top} mostró dinámicamente la lista de procesos activos, ordenados por
defecto según el consumo de CPU, y actualizó la información cada tres segundos. En la cabecera
se presentaron métricas globales del sistema: tiempo de actividad (\emph{uptime}), número de
usuarios conectados, carga promedio en 1, 5 y 15 minutos, estadísticas de tareas y el uso de
CPU y memoria RAM. La sesión se terminó mediante la tecla \texttt{q}.

\begin{lstlisting}[language=bash,caption={Ejecución de la herramienta top.},label={lst:top}]
$ top
\end{lstlisting}

\begin{figure}[H]
    \centering
    \includegraphics[width=0.9\textwidth]{figura7_top.png}
    \caption[Interfaz de top.]{Interfaz de la herramienta \texttt{top} mostrando procesos ordenados por consumo de CPU.}
    \label{fig:top}
\end{figure}

\subsubsection{\texttt{htop} — Monitor interactivo mejorado}

La herramienta \texttt{htop} constituyó una versión extendida e interactiva de \texttt{top}, que presentó
la misma información con una interfaz gráfica en color, barras de progreso para CPU y RAM y
soporte para navegación con el ratón. Permitió ordenar la lista de procesos de forma interactiva
y gestionar procesos directamente desde la interfaz sin necesidad de conocer el PID.

\begin{lstlisting}[language=bash,caption={Ejecución de la herramienta htop.},label={lst:htop}]
$ htop
\end{lstlisting}

\begin{figure}[H]
    \centering
    \includegraphics[width=0.9\textwidth]{figura8_htop.png}
    \caption[Interfaz de htop.]{Interfaz de \texttt{htop} con barras de progreso por núcleo y métricas de memoria.}
    \label{fig:htop}
\end{figure}

\subsection{Árbol de procesos con \texttt{pstree}}

La herramienta \texttt{pstree} representó la jerarquía de todos los procesos del sistema en
forma de árbol, partiendo del proceso \texttt{init} (PID~1) o \texttt{systemd} como raíz.

\subsubsection{\texttt{pstree} — Árbol básico}

\begin{lstlisting}[language=bash,caption={Visualización básica del árbol de procesos con pstree.},label={lst:pstree}]
$ pstree
\end{lstlisting}

\begin{figure}[H]
    \centering
    \includegraphics[width=0.85\textwidth]{figura9_pstree.png}
    \caption{Árbol de procesos básico generado con \texttt{pstree}.}
    \label{fig:pstree}
\end{figure}

\subsubsection{\texttt{pstree -p} — Con identificadores de proceso}

La opción \texttt{-p} añadió el PID de cada proceso entre paréntesis junto a su nombre:

\begin{lstlisting}[language=bash,caption={Árbol de procesos con PID mediante pstree -p.},label={lst:pstreep}]
$ pstree -p
\end{lstlisting}

\begin{figure}[H]
    \centering
    \includegraphics[width=0.85\textwidth]{figura10_pstree_p.png}
    \caption[Árbol de procesos con PID.]{Árbol de procesos con PID visible mediante \texttt{pstree -p}.}
    \label{fig:pstreep}
\end{figure}

\subsubsection{\texttt{pstree -u root} — Filtrado por usuario root}

La opción \texttt{-u} mostró el nombre del usuario propietario de cada proceso cuando este
cambió respecto al del proceso padre, lo que permitió identificar transiciones de privilegios:

\begin{lstlisting}[language=bash,caption={Árbol de procesos del usuario root con pstree -u root.},label={lst:ptreeu}]
$ pstree -u root
\end{lstlisting}

\begin{figure}[H]
    \centering
    \includegraphics[width=0.85\textwidth]{figura11_pstree_u.png}
    \caption[Árbol de procesos del usuario root.]{Árbol de procesos del usuario \texttt{root} con \texttt{pstree -u root}.}
    \label{fig:ptreeu}
\end{figure}

\subsubsection{\texttt{pstree -a} — Con argumentos completos}

La opción \texttt{-a} mostró, junto al nombre de cada proceso, la línea de comandos completa
con la que fue invocado, así como todos sus argumentos:

\begin{lstlisting}[language=bash,caption={Árbol de procesos con argumentos completos mediante pstree -a.},label={lst:ptreea}]
$ pstree -a
\end{lstlisting}

\begin{figure}[H]
    \centering
    \includegraphics[width=0.85\textwidth]{figura12_pstree_a.png}
    \caption[Árbol con argumentos completos.]{Árbol de procesos con argumentos completos mediante \texttt{pstree -a}.}
    \label{fig:ptreea}
\end{figure}

\subsection{Búsqueda de procesos: \texttt{pgrep}, \texttt{pidof} y \texttt{sleep}}

\subsubsection{\texttt{pgrep} — Búsqueda por nombre o patrón}

El comando \texttt{pgrep} (\emph{Process Grep}) buscó, entre los procesos en ejecución, aquellos
cuyo nombre coincidió con el patrón proporcionado y devolvió sus PIDs. A diferencia de canalizar
la salida de \texttt{ps} con \texttt{grep}, \texttt{pgrep} actuó directamente sobre la tabla de
procesos del kernel:

\begin{lstlisting}[language=bash,caption={Búsqueda de procesos por nombre con pgrep.},label={lst:pgrep}]
$ pgrep firefox
1234
\end{lstlisting}

\subsubsection{\texttt{pidof} — PID por nombre exacto}

El comando \texttt{pidof} devolvió el o los PIDs de todos los procesos cuyo nombre de ejecutable
coincidió exactamente con el argumento especificado:

\begin{lstlisting}[language=bash,caption={Obtención del PID por nombre exacto con pidof.},label={lst:pidof}]
$ pidof firefox
1234 5678
\end{lstlisting}

\subsubsection{\texttt{sleep} — Proceso de espera controlada}

El comando \texttt{sleep} creó un proceso que se mantuvo en estado \emph{Bloqueado} durante
el número de segundos indicado. Se empleó para generar procesos de prueba cuyo PID pudiera
ser identificado con los comandos anteriores:

\begin{lstlisting}[language=bash,caption={Uso del comando sleep para crear procesos de prueba.},label={lst:sleep}]
$ sleep 5        # suspende 5 segundos
$ sleep 1000     # suspende 1000 segundos (proceso de prueba)
$ pidof sleep    # devuelve el PID del proceso sleep en ejecucion
\end{lstlisting}

\subsection{Sistema de archivos \texttt{/proc}}

\subsubsection{\texttt{ls /proc} — Listado del sistema de archivos virtual}

\begin{lstlisting}[language=bash,caption={Listado del contenido del sistema de archivos virtual /proc.},label={lst:lsproc}]
$ ls /proc
\end{lstlisting}

\begin{figure}[H]
    \centering
    \includegraphics[width=0.85\textwidth]{figura13_ls_proc.png}
    \caption[Contenido de /proc.]{Contenido del sistema de archivos virtual \texttt{/proc}.}
    \label{fig:lsproc}
\end{figure}

Cada directorio numérico correspondió al PID de un proceso activo; los archivos y directorios
con nombres textuales contuvieron información global del sistema.

\subsubsection{\texttt{cat /proc/cpuinfo} — Información del procesador}

El archivo \texttt{/proc/cpuinfo} expuso en texto plano la información detallada de cada
núcleo del procesador: modelo, velocidad en MHz, tamaño de caché y banderas de instrucciones
soportadas, entre otros:

\begin{lstlisting}[language=bash,caption={Lectura de la información del procesador desde /proc/cpuinfo.},label={lst:cpuinfo}]
$ cat /proc/cpuinfo
\end{lstlisting}

\begin{figure}[H]
    \centering
    \includegraphics[width=0.85\textwidth]{figura14_cpuinfo.png}
    \caption[Información del procesador desde /proc/cpuinfo.]{Información del procesador obtenida desde \texttt{/proc/cpuinfo}.}
    \label{fig:cpuinfo}
\end{figure}

\subsubsection{\texttt{lscpu} — Resumen de la arquitectura del procesador}

La herramienta \texttt{lscpu} presentó de forma estructurada la arquitectura del procesador,
el número de CPUs, de núcleos, de hilos por núcleo, la velocidad máxima y mínima en MHz y
la información de caché:

\begin{lstlisting}[language=bash,caption={Resumen de la arquitectura del procesador con lscpu.},label={lst:lscpu}]
$ lscpu
\end{lstlisting}

\begin{figure}[H]
    \centering
    \includegraphics[width=0.85\textwidth]{figura15_lscpu.png}
    \caption{Arquitectura del procesador obtenida con \texttt{lscpu}.}
    \label{fig:lscpu}
\end{figure}

\subsection{Monitoreo de la CPU con \texttt{mpstat}}

La herramienta \texttt{mpstat} (\emph{Multiprocessor Statistics}) formó parte del paquete
\texttt{sysstat} y reportó estadísticas de uso del procesador.

\subsubsection{\texttt{mpstat} — Estadísticas globales del procesador}

\begin{lstlisting}[language=bash,caption={Estadísticas globales del procesador con mpstat.},label={lst:mpstat}]
$ mpstat
\end{lstlisting}

\begin{figure}[H]
    \centering
    \includegraphics[width=0.85\textwidth]{figura16_mpstat.png}
    \caption{Estadísticas globales del procesador obtenidas con \texttt{mpstat}.}
    \label{fig:mpstat}
\end{figure}

Los campos principales incluyeron \texttt{\%usr} (tiempo en modo usuario), \texttt{\%sys}
(tiempo en modo kernel), \texttt{\%iowait} (tiempo esperando E/S) e \texttt{\%idle}
(tiempo de inactividad).

\subsubsection{\texttt{mpstat -P all 1} — Estadísticas por núcleo}

La opción \texttt{-P all} reportó el uso de cada núcleo del procesador de forma individual;
el argumento \texttt{1} indicó actualización cada segundo:

\begin{lstlisting}[language=bash,caption={Estadísticas por núcleo del procesador con mpstat -P all 1.},label={lst:mpstatall}]
$ mpstat -P all 1
\end{lstlisting}

\begin{figure}[H]
    \centering
    \includegraphics[width=0.85\textwidth]{figura17_mpstat_all.png}
    \caption[Estadísticas por núcleo con mpstat.]{Uso del procesador por núcleo con \texttt{mpstat -P all 1}.}
    \label{fig:mpstatall}
\end{figure}

% ==================================================================
\section{Conclusiones}
% ==================================================================

El estudio y la aplicación práctica de las herramientas de gestión y monitoreo de procesos en
Linux permitieron al estudiante desarrollar competencias esenciales para la administración de
sistemas operativos. El comando \texttt{ps} y sus variantes demostraron ser la base fundamental
para la consulta estática del estado de los procesos: \texttt{ps -aux} proporcionó una visión
global de todos los procesos del sistema, \texttt{ps -ef} expuso las relaciones padre-hijo
mediante la columna PPID, \texttt{ps -u root} permitió auditar los procesos con privilegios
administrativos y \texttt{ps -e --forest} visualizó la jerarquía completa en formato de árbol.

Las herramientas \texttt{pgrep}, \texttt{pidof} y \texttt{sleep} facilitaron la identificación
y creación de procesos de prueba de forma controlada, complementando la gestión dinámica de
la terminal. La exploración del sistema de archivos virtual \texttt{/proc} reveló la arquitectura
interna con la que el kernel de Linux expone la información de los procesos al espacio de usuario,
que constituyó la fuente de datos primaria de todas las herramientas de monitoreo estudiadas.

Las herramientas de monitoreo dinámico revisadas —\texttt{top}, \texttt{htop}, \texttt{pstree}
y \texttt{mpstat}— presentaron distintos niveles de profundidad y especialización. Para
diagnósticos rápidos e interactivos en sistemas de producción resultó suficiente \texttt{htop},
dada su interfaz ncurses con barras de progreso por núcleo y soporte para el envío de señales
desde la propia pantalla; mientras que \texttt{mpstat} constituyó la opción más apropiada para
el análisis cuantitativo del uso del procesador por núcleo físico y lógico.

% ==================================================================
% Referencias
% ==================================================================

\begin{thebibliography}{99}
\addcontentsline{toc}{section}{Referencias}
\setlength{\itemsep}{0.4em}

\bibitem{kerrisk2010} Kerrisk, M. (2010). \emph{The Linux programming interface: A Linux and
UNIX system programming handbook}. No Starch Press.

\bibitem{love2010} Love, R. (2010). \emph{Linux kernel development} (3.ª~ed.). Addison-Wesley.

\bibitem{silberschatz2018} Silberschatz, A., Galvin, P. B., y Gagne, G. (2018). \emph{Operating
system concepts} (10.ª~ed.). John Wiley \& Sons.

\bibitem{tanenbaum2015} Tanenbaum, A. S., y Bos, H. (2015). \emph{Modern operating systems}
(4.ª~ed.). Pearson Education.

\bibitem{linuxfoundation2024} The Linux Foundation. (2024). \emph{The Linux man pages project}.
Recuperado de \url{https://man7.org/linux/man-pages/}

\end{thebibliography}

% ==================================================================
\end{document}
% ==================================================================
